\chapter{Introducción}
\label{ch:introduccion}
% Redacción impersonal; sin notas al pie ni referencias en este capítulo (recomendación ETSINF).

En poco más de una década, el aprendizaje automático ha pasado de ser una disciplina de interés sobre todo académico a sostener buena parte de la tecnología de uso cotidiano. Los sistemas que reconocen objetos en imágenes, transcriben voz o generan texto comparten un mismo fundamento, redes neuronales profundas entrenadas sobre grandes volúmenes de datos. Buena parte de los avances del periodo tiene además un origen común, la escala, porque más parámetros, más datos y más tiempo de entrenamiento suelen bastar para mejorar los resultados. La receta funciona y su éxito ha desplazado el cuello de botella de la disciplina. Hoy lo que limita la investigación rara vez son las ideas o los datos. Es el cómputo disponible para ponerlos a prueba.

El coste de esa escala recae primero sobre la industria. Los laboratorios que entrenan los modelos más grandes planifican el trabajo en meses de GPU, con la factura eléctrica que eso supone, y la capacidad de cálculo decide qué organizaciones pueden abordar qué problemas. El coste no es solo económico. El consumo energético del entrenamiento ha hecho de la eficiencia un objeto de estudio, con una línea de trabajo asociada, el llamado \emph{green AI}, que propone valorar los avances también por lo que cuestan.

En la universidad la misma restricción se vive desde el otro lado. Un grupo con menos medios dispone de unas pocas GPU y de un número cerrado de horas de cálculo, y la pregunta pasa de qué modelo entrenar a cómo repartir las horas entre las configuraciones que merezca la pena probar.

Industria y universidad llegan así a la misma pregunta práctica, cuáles de las configuraciones que compiten por las mismas horas merecen seguir y cuánto hay que esperar para saberlo. Esta memoria estudia una de las señales que podrían responderla, los gradientes de la red mientras aprende, y la pone a prueba donde esa decisión se toma, entre entrenamientos de una misma red que solo difieren en el \emph{learning rate}.

Este capítulo presenta el problema que motiva el trabajo y enmarca lo que sigue. La Sección~\ref{sec:motivacion} expone por qué interesa predecir la eficiencia de un entrenamiento desde sus primeras \emph{epochs} y dónde está el hueco que esta memoria pretende cubrir. La Sección~\ref{sec:objetivos} formula la pregunta de investigación y los objetivos que de ella se derivan. La Sección~\ref{sec:alcance} delimita el estudio y separa lo que entra de lo que queda fuera. La Sección~\ref{sec:estructura} cierra el capítulo con un recorrido por el resto de la memoria.

\section{Motivación}
\label{sec:motivacion}

Cada vez que se elige una arquitectura, un optimizador o un conjunto de hiperparámetros, la forma habitual de comprobar si la elección es buena es lanzar el entrenamiento completo y medir el resultado. Se repite muchas veces, porque el espacio de configuraciones posibles es enorme y rara vez se sabe de antemano cuál funcionará mejor.

Buena parte de ese coste se consume en entrenamientos que podrían descartarse pronto. Existen criterios cuantitativos para cortarlos, y todos leen la curva de validación o la de \emph{loss}: extrapolar la curva de aprendizaje desde sus primeras \emph{epochs}, o repartir el presupuesto entre muchas configuraciones y eliminar las peores por rondas, como hace Hyperband. Esas curvas son la señal que ya se mira. La cuestión que plantea este trabajo es si el gradiente contiene información que esas curvas no contienen. Si la contiene, las primeras \emph{epochs} de un presupuesto fijado de antemano bastarían para decidir mejor entre configuraciones y reservar el cómputo para las más prometedoras.

Entre las señales disponibles durante el entrenamiento de una red neuronal, el gradiente es una de las más naturales para observar el progreso del aprendizaje, porque es la información que el optimizador utiliza para avanzar en cada paso. La curva de \emph{loss} resume ese avance en un número por paso. El gradiente de cada ejemplo conserva una estructura que ese promedio descarta, hacia dónde empuja cada ejemplo y cuánto se separan unos de otros. Dos aspectos de esa estructura han recibido atención en trabajos distintos.

El primero es su alineación direccional, esto es, hacia dónde apuntan los gradientes de ejemplos distintos, si al mismo lado o en sentidos opuestos. Es una cuestión de dirección, no de magnitud, que se resume en el ángulo entre ellos. Cuando distintos ejemplos empujan en la misma dirección, sus gradientes se refuerzan y el paso de descenso beneficia a muchos a la vez, mientras que cuando se oponen se cancelan entre sí y el avance se estanca.

El segundo es su variabilidad estocástica. El optimizador ve un \emph{minibatch} cada vez y no todos los datos a la vez, así que el gradiente con el que avanza no es el gradiente verdadero sobre el conjunto completo, sino una estimación a partir de una muestra, y esa estimación cambia según qué ejemplos toquen. La variabilidad mide cuánto se dispersan esos gradientes alrededor de su media, la distancia de cada uno al promedio. Cuando la dispersión es grande, el optimizador recibe estimaciones ruidosas y avanza con dificultad.

En los últimos años se han propuesto varias métricas para cuantificar estos dos aspectos. Cada una se ha presentado en su propio trabajo, con su propio conjunto de datos y su propia arquitectura, y casi siempre con un objetivo distinto, como explicar por qué una red generaliza, elegir el tamaño de \emph{batch} o decidir cuándo detener el entrenamiento. Tres piezas de la pregunta de esta memoria tienen antecedentes. Comparar decenas de medidas sobre un mismo banco de modelos ya se ha hecho para la generalización. Cuán temprano basta medir es la pregunta de la extrapolación de curvas de aprendizaje. Si una medida del gradiente aporta algo sobre una referencia trivial es la pregunta que la búsqueda de arquitecturas se hizo con los \emph{proxies} de coste cero, medidos antes de entrenar, y la respuesta fue muchas veces que no. Lo que sigue sin cubrir es si una métrica del gradiente, medida en las primeras \emph{epochs}, ordena las configuraciones de una misma arquitectura que solo difieren en el \emph{learning rate} mejor que la curva de validación de esas mismas \emph{epochs}, bajo SGD y bajo Adam, y qué cuesta decidir con ella. El Capítulo~\ref{ch:estado-del-arte} sitúa cada antecedente.

Este trabajo cubre ese hueco con un estudio correlacional sistemático. Se entrena una matriz amplia de configuraciones y arquitecturas de redes neuronales para distintos problemas de visión por computador, se calculan en todas ellas las métricas a lo largo del entrenamiento, y se cuantifica en qué medida cada métrica, leída en varias ventanas del proceso, anticipa la eficiencia del entrenamiento completo. El criterio de éxito no es solo que una métrica se asocie con la eficiencia, sino que supere a un predictor de referencia (\emph{baseline}) nombrado de antemano, la curva de validación leída en esa misma ventana, y que no sea esa referencia con otro nombre. El coste de cada métrica se da en su clase de complejidad y en el sobrecoste medido sobre la matriz. Un resultado negativo bien establecido, como por ejemplo la evidencia de que esas métricas no aportan nada que la curva de validación no contenga ya, sería también una contribución útil, porque ahorraría a otros el coste de calcularlas.

\section{Objetivos}
\label{sec:objetivos}

La pregunta de investigación que vertebra la memoria es la siguiente:

\begin{quote}
\itshape ¿Puede alguna métrica del gradiente, medida en las primeras \emph{epochs}, predecir la eficiencia del entrenamiento completo?
\end{quote}

En esta memoria la eficiencia del entrenamiento reúne tres aspectos. La \emph{velocidad de convergencia} es el número de \emph{epochs} que cuesta alcanzar el nivel propio de esa arquitectura en ese problema, leído sobre la curva de validación. El \emph{rendimiento final} es el \emph{accuracy} del modelo sobre el conjunto de test, que se evalúa una sola vez al terminar. La \emph{generalización} es la diferencia entre lo que el modelo logra sobre datos de entrenamiento y sobre datos que no ha visto. Solo el primero es eficiencia en el sentido de rendimiento por unidad de cómputo, y la palabra se usa aquí en ese sentido amplio. Cada aspecto se concreta en una o varias variables en el Capítulo~\ref{ch:metodologia}, y la pregunta se evalúa contra los tres.

La pregunta se responde donde el diseño puede responderla. Con el conjunto de datos, la arquitectura y el optimizador fijos, lo único que cambia entre entrenamientos es el \emph{learning rate} y la \emph{seed}, y toda correlación se calcula dentro de esa celda. Lo que el estudio mide es, por tanto, si una métrica temprana ordena los \emph{learning rates} de una celda mejor que la validación temprana, y cuánto \emph{accuracy} cuesta elegir el \emph{learning rate} con ella. El caso de uso es elegir entre configuraciones, no leer el nivel absoluto de la métrica en un entrenamiento aislado.

El objetivo general del trabajo es cuantificar en qué medida las métricas tempranas de gradiente predicen la eficiencia del entrenamiento completo en tareas de clasificación de imágenes, e identificar cuáles ofrecen la mejor capacidad predictiva y a qué coste. Se concreta en seis objetivos específicos. Los dos primeros son el núcleo del trabajo, y los otros cuatro leen el mismo resultado desde otro ángulo:

\begin{enumerate}
  \item \textbf{OE1 (existencia).} Determinar si al menos una métrica temprana de gradiente se relaciona de forma consistente con la eficiencia en las condiciones estudiadas. Alguna asociación es esperable, porque dentro de una celda el \emph{learning rate} mueve a la vez la métrica y la eficiencia. Por eso este objetivo es un requisito previo de los otros cinco, que sin él carecen de objeto. \label{obj:existencia}
  \item \textbf{OE2 (valor incremental).} Comprobar si alguna métrica supera a una referencia que no cuesta nada, la curva de validación leída en esa misma ventana, sin ser esa referencia con otro nombre. Se decide sobre el rendimiento final y la generalización, porque la velocidad se lee de la misma curva que la referencia. Para el rendimiento se añade una lectura práctica, cuánto \emph{accuracy} se pierde al elegir el \emph{learning rate} con la métrica en vez de con la validación. \label{obj:incremental}
  \item \textbf{OE3 (comparación de familias).} Ordenar las métricas por su capacidad predictiva y ver si las primeras posiciones las ocupa una sola familia. Las dos familias comparten parte de lo que miden, así que la conclusión sobre familias es cualitativa. \label{obj:familias}
  \item \textbf{OE4 (suficiencia temprana).} Determinar cuán temprano basta medir, comparando las ventanas del 5 al 50\,\% del presupuesto. Para la velocidad solo sirven las dos primeras, por una razón que el Capítulo~\ref{ch:metodologia} expone. \label{obj:temprana}
  \item \textbf{OE5 (robustez entre optimizadores).} Evaluar si el sentido de la relación se preserva al cambiar de optimizador, de descenso por gradiente estocástico (SGD) a Adam, comparando cada celda con la que solo difiere en el optimizador. \label{obj:robustez}
  \item \textbf{OE6 (concordancia con la literatura).} Comprobar si el sentido de cada relación observada coincide con el que predice el trabajo que propuso la métrica. Es la lectura con signo de OE1, y solo se contrasta donde el artículo enuncia una predicción; las demás predicciones son extrapolaciones de esta memoria y se reportan sin contrastar. \label{obj:concordancia}
\end{enumerate}

Estos seis objetivos se formalizan en el Capítulo~\ref{ch:metodologia} como las hipótesis H1 a H6, en correspondencia uno a uno con OE1 a OE6, y cada una se enuncia junto al criterio que bastaría para refutarla. El Capítulo~\ref{ch:conclusiones} retoma la lista para examinar, uno a uno, en qué medida se han alcanzado.

\section{Alcance y limitaciones}
\label{sec:alcance}

El trabajo es un estudio \emph{correlacional}. Su finalidad es medir la asociación entre las métricas tempranas y la eficiencia, no intervenir en el entrenamiento a partir de esa señal. El trabajo de diseñar un sistema de intervención, por ejemplo detener el entrenamiento antes de tiempo o ajustar el \emph{learning rate} en función de la alineación medida, queda fuera del alcance y se recoge como línea de trabajo futuro en el Capítulo~\ref{ch:trabajos-futuros}. Como en todo estudio correlacional, las relaciones que se encuentren no deben interpretarse como causales.

El estudio se desarrolla en visión por computador, en concreto en clasificación de imágenes. Abarca cuatro conjuntos de datos de dificultad creciente, tres arquitecturas (una red \emph{fully connected}, una red convolucional simple y una ResNet-18) y dos optimizadores (SGD y Adam). La combinación da lugar a una matriz de experimentos que, con ocho \emph{learning rates} y cinco \emph{seeds} por configuración, suma 960 entrenamientos. El detalle de esta matriz se expone en el Capítulo~\ref{ch:metodologia}.

Dentro de cada celda solo se barre el \emph{learning rate}, con cinco \emph{seeds} por valor. Todo lo demás, del tamaño de \emph{batch} al presupuesto de \emph{epochs}, es idéntico en los 960 entrenamientos y se detalla en el Capítulo~\ref{ch:metodologia}. La configuración de entrenamiento es deliberadamente simple, con \emph{learning rate} constante y sin regularización ni \emph{data augmentation}, y el precio es que los valores de \emph{accuracy} quedan por debajo de los publicados. Las conclusiones valen para esas condiciones.

El método de análisis se fijó después de ejecutar la matriz y antes de calcular ningún coeficiente entre una métrica y una variable de eficiencia, y se revisó después en diez puntos. La Sección~\ref{sec:protocolo-cuando} da las fechas y los motivos.

Quedan fuera del alcance los dominios distintos de la visión, las arquitecturas de tipo \emph{transformer} y los conjuntos de datos a gran escala. El conjunto ImageNet completo se sustituye por una versión reducida, Tiny-ImageNet, por restricciones de disco y de cómputo, y el Capítulo~\ref{ch:metodologia} declara las consecuencias.

\section{Estructura de la memoria}
\label{sec:estructura}

El resto de la memoria se organiza como sigue. El Capítulo~\ref{ch:estado-del-arte} revisa críticamente la literatura, ordenándola en torno a las dos familias de métricas y a los predictores de eficiencia que sirven de referencia, y termina situando la aportación de este trabajo frente a lo existente. El Capítulo~\ref{ch:fundamentos} fija las definiciones formales que el estado del arte deja enunciadas: el descenso por gradiente estocástico, cada una de las métricas y los indicadores de eficiencia. El Capítulo~\ref{ch:metodologia} presenta el diseño experimental: las hipótesis con sus criterios falsables, las variables, la ventana temporal, la configuración del entrenamiento y la matriz de experimentos, la calibración de presupuestos y umbrales, los protocolos de evaluación y de análisis, y los riesgos que el diseño debe neutralizar. El Capítulo~\ref{ch:desarrollo} describe el desarrollo del software que hace posible el estudio. El Capítulo~\ref{ch:resultados} reúne los resultados de los experimentos. El Capítulo~\ref{ch:conclusiones} interpreta esos resultados y revisa el cumplimiento de los objetivos, y el Capítulo~\ref{ch:trabajos-futuros} recoge las líneas de continuación que quedan abiertas. Cierran la memoria los anexos, que recogen la reflexión sobre los Objetivos de Desarrollo Sostenible y los detalles de configuración necesarios para reproducir el estudio.
